\section{Hipoteza warunkowego przyspieszenia}

Hipotezę pracy można sformułować jako hipotezę warunkowego przyspieszenia automatyzacji. W ujęciu intuicyjnym brzmi ona następująco: instrument zaprojektowany jako osłona społeczna przed skutkami automatyzacji może sam zmienić ceny względne i warunki finansowania w taki sposób, że stanie się bodźcem do automatyzacji. Odwróceniu ulega zwyczajowy porządek przyczynowy. BDP nie jest wyłącznie skutkiem rozwoju AI; może być jednym z mechanizmów, które przyspieszają decyzje firm o zastępowaniu pracy kapitałem.

Ta hipoteza ma jednak dwie wersje. Wersja mocna jest monotoniczna: im wyższy BDP, tym silniejsza presja automatyzacyjna i tym większa adopcja AI/hybrid. Wersja słabsza, ale ekonomicznie bardziej realistyczna, jest warunkowa: deficytowo finansowany BDP może przyspieszać automatyzację tylko do momentu, w którym sektor prywatny zachowuje zdolność finansowania inwestycji. Po przekroczeniu pewnego progu ten sam transfer może ograniczać automatyzację nie dlatego, że sam wzrost długu mechanicznie wypiera nakłady inwestycyjne, lecz dlatego, że presja płacowo-rezerwowa tworzy potrzebę automatyzacji, a pogorszenie warunków fiskalno-finansowych zmniejsza przestrzeń bilansową firm do sfinansowania tej reakcji. Zakres tej pracy dotyczy właśnie wariantu deficytowego; finansowanie przez VAT, PIT, cięcia wydatków lub warianty mieszane mogłoby zmienić kanały transmisji.

Mikroekonomiczny komponent hipotezy opiera się na płacy rezerwowej. Jeżeli dochód niezależny od pracy rośnie, część gospodarstw domowych może wymagać wyższej płacy, aby zaakceptować zatrudnienie. Literatura teoretyczna i empiryczna dotycząca podatków, transferów oraz dochodu podstawowego nie wskazuje na jedną stabilną elastyczność podaży pracy; efekty zależą od konstrukcji programu, populacji i otoczenia rynku pracy \parencite{mirrlees1971,saez2002,vanParijs2017,forget2011,marinescu2018,kangas2019,jonesMarinescu2022,hoynesRothstein2019,banerjee2019}. Dlatego \(\lambda = 0{,}5\) należy traktować jako parametr scenariuszowy po stronie relatywnie silnego przełożenia transferu na płacę rezerwową, a nie jako bezpośrednią estymację empiryczną. W modelu \amor{} kanał ten jest zapisany jako:
\begin{equation}
  w^{res}_t = w^{min}_t + \lambda \cdot BDP,
\end{equation}
gdzie \(w^{res}_t\) oznacza płacę rezerwową, \(w^{min}_t\) płacę minimalną, a \(\lambda\) parametr przełożenia transferu na płacę rezerwową. W wariancie centralnym przyjęto \(\lambda = 0{,}5\). Oznacza to, że gospodarstwo domowe otrzymuje pełny transfer BDP, natomiast tylko część tego transferu wpływa w modelu na płacę rezerwową. Parametr \(\lambda\) nie zmniejsza samego transferu; opisuje jedynie siłę kanału negocjacyjno-płacowego. Ponieważ jest to założenie behawioralne, analiza odporności \(\lambda\) jest częścią głównego wyniku, a nie dodatkiem technicznym.

Makroekonomiczny komponent hipotezy przebiega przez popyt, inflację, politykę pieniężną, dług publiczny, kurs walutowy i koszt finansowania. Transfer zwiększa dochód gospodarstw domowych, lecz nie każdy dodatkowy złoty musi stać się krajowym popytem finalnym. Część impulsu może trafić do oszczędności, importu albo bilansów gospodarstw. Jednocześnie finansowanie BDP zwiększa wydatki publiczne, deficyt i dług. W gospodarce otwartej, takiej jak Polska, dodatkowym kanałem jest kurs PLN/EUR: wzrost popytu i importochłonności może pogarszać rachunek bieżący, osłabiać walutę i zwiększać ceny importowanych komponentów oraz technologii.

Dla firmy najważniejsza jest kombinacja tych kanałów. Jeśli BDP podnosi koszt pracy, bodziec do automatyzacji rośnie; jeśli równocześnie rośnie popyt, firma może mieć większą motywację do zwiększenia mocy produkcyjnych. Jeżeli jednak rośnie koszt długu publicznego, rentowność obligacji, koszt długu korporacyjnego albo ryzyko kursowe, finansowanie nakładów inwestycyjnych staje się trudniejsze. Tak powstaje próg absorpcji: obszar, w którym pozytywny impuls popytowy i presja kosztu pracy przestają dominować nad wpływem kosztu kapitału, długu publicznego, kursu walutowego i słabszych inwestycji prywatnych. Nie jest to próg długu sam w sobie; kluczowa jest interakcja bodźca do automatyzacji z możliwością jej sfinansowania.

Hipoteza jest zatem falsyfikowalna. Monotoniczny wzrost adopcji AI/hybrid potwierdzałby wersję mocną; brak reakcji albo spadek od początku odrzucałby tezę o przyspieszającej roli BDP; lokalny wzrost przy umiarkowanym transferze i spadek po przekroczeniu wyższego poziomu BDP wspierałby wersję warunkową. Dlatego kluczowe nie jest samo stwierdzenie, że BDP „przyspiesza” albo „nie przyspiesza” automatyzacji, lecz ustalenie, czy działa jako katalizator tylko w określonym korytarzu makro-finansowym.
